\section{Discussion}\label{sec:discussion}

\subsection{What was reproduced, improved, or not reproduced}
\TODO{verdict table: reproduced / improved / not reproduced -- after GNN runs}

\subsection{Threats to validity}\label{sec:threats}
\paragraph{Internal validity (does our pipeline do what the paper did?).}
\begin{itemize}
  \item \emph{Unreleased implementation.} ICFG construction above LChecker's CFG builder, the pruning rules, the
        built-in function list, the risk lists and the statement text format are our reading of the paper. Different
        choices change the vocabulary (we obtain 484k distinct normalised statements) and therefore the embeddings.
  \item \emph{Extractor fixes.} We changed LChecker's behaviour on namespaces, try/catch, \texttt{for} conditions and
        short tags (\Cref{sec:method}). Each fix was needed to avoid graphs that are empty for one class only; if the
        authors ran vanilla LChecker on PHP~8 their graphs were partly empty, and their numbers are not directly
        comparable to ours in either direction.
  \item \emph{Unspecified training details} (epochs for Table 2, early stopping, model selection, embedding training
        corpus, threshold). We fixed them before looking at test results and report all of them (\Cref{tab:hyper}).
  \item \emph{Resource-driven deviations}: 2,000-statement cap (15\% of webshells truncated), FastText 200k buckets,
        own GloVe implementation.
\end{itemize}

\paragraph{Construct validity (do the labels and metrics measure webshell detection?).}
\begin{itemize}
  \item \emph{Source-level labels.} ``Benign'' means ``file of a large, modern framework/CMS release'', ``webshell'' means
        ``file of a public collection''. The two populations also differ in age (PHP~4/5 vs.\ PHP~8 style), size,
        structure (procedural vs.\ namespaced classes) and parseability (\Cref{tab:extraction}). A classifier can reach
        high scores by separating these \emph{populations} rather than by recognising malicious behaviour; the
        graph-statistics probe and the hand-written standalone benign scripts (\Cref{sec:errors}) quantify this.
  \item \emph{Label noise.} Collections contain non-malicious helper files; benign projects contain test fixtures and
        installers with webshell-like capabilities. We did not relabel by hand.
  \item \emph{Class prior.} 23\% webshells in test is unrealistically high for a file-upload monitor, so precision and F1
        overstate deployment performance; FPR on the 9,091 held-out benign files is the more relevant number.
\end{itemize}

\paragraph{External validity (does it generalise?).}
\begin{itemize}
  \item Only five benign projects in training; four more in testing. Industrial web servers (the paper's motivating
        setting) run vendor-specific, often small, procedural PHP applications that resemble neither population.
  \item Only 18 webshells are novel w.r.t.\ the paper sources; recall on unseen families is therefore estimated with
        very wide intervals. No temporal split was possible (\Cref{sec:leak}).
  \item No adaptive attacker: renaming variables and functions is neutralised by normalisation, but string building,
        dynamic calls and payload encoding already defeat the risk mechanism for many collected samples.
\end{itemize}

\paragraph{Conclusion validity (are the differences real?).}
Five seeds give 95\% intervals of roughly $\pm$0.01--0.03 F1 for the GNNs; several differences reported in the paper
(e.g.\ between embedding methods, 0.91--0.95 F1 in a single run) are of the same order and cannot be distinguished
without repeated runs. We never select hyper-parameters on the test set.
